% Lösungen Einheit 1 von 3 – Pyramide (zusammenbau v0.1)
\einheitenkopf{Einheit 1 von 3 · Pyramide}

\erg{12}{a) Spitze; Grundfläche: das Quadrat unten}
\erg{13}{a) Höhe von der Spitze senkrecht zur Mitte der Grundfläche, von dort zur Mitte einer Grundkante (halbe Grundkante, $3$ m), zurück zur Spitze; die Seitenhöhe ist die Hypotenuse}
\erg{14}{a) Oberfläche – Stoff ist Hülle}
\erg{15}{a) $V = \frac{1}{3} \cdot G \cdot h$ mit $G = 30$ cm², $h = 8$ cm}
\erg{16}{a) Höhe \quad b) $V = 27 \cdot 8 : 3 = 72$ cm³ \quad c) $G = 6 \cdot 6 = 36$ cm²; $V = 36 \cdot 7 : 3 = 84$ cm³ \quad d) $G = 8 \cdot 5 = 40$ cm²; $V = 40 \cdot 9 : 3 = 120$ cm³ \quad e) $G = 3{,}8^2 = 14{,}44$ cm²; $V = 14{,}44 \cdot 5{,}5 : 3 \approx 26{,}5$ cm³ \quad f) halbe Grundkante $8$ cm; $h_s = \sqrt{15^2 + 8^2} = \sqrt{289} = 17$ cm \quad g) $A = 8 \cdot 11 : 2 = 44$ cm² \quad h) $M = 4 \cdot 6 \cdot 9 : 2 = 108$ cm² \quad i) $G = 64$ cm², $M = 4 \cdot 8 \cdot 7 : 2 = 112$ cm², $O = 176$ cm² \quad j) (halbe Diagonale)$^2 = 3^2 + 3^2 = 18$; Seitenkante $= \sqrt{7^2 + 18} = \sqrt{67} \approx 8{,}19$ cm \quad k) alle vier Seiten des Quadrats $4$ cm; in jedem Dreieck die Seitenhöhe von der Spitze zur Mitte der Quadratseite, je $5$ cm \quad l) Grundquadrat: vordere Kante $4$ cm, Tiefenkanten $2$ cm unter 45°; Diagonalen schneiden sich im Fußpunkt; Höhe $5$ cm senkrecht darüber, gestrichelt; Spitze mit den vier Ecken verbinden \quad m) $h = 3 \cdot 96 : 36 = 8$ cm \quad n) $M = 4 \cdot 5 \cdot 3{,}2 : 2 = 32$ m² \quad o) Diagonalen der Grundfläche einzeichnen; vom Schnittpunkt senkrecht zur Spitze gestrichelt, mit $1{,}8$ m beschriften; eine Grundkante mit $4{,}5$ m beschriften}

16k)

\netzpyramide{2}{2.5}{$4$ cm}{$5$ cm}


16l)

\pyramide{4}{4}{5}{$4$ cm}{}{$5$ cm}

\erg{17}{a) Jonas hat den Faktor $\frac{1}{3}$ vergessen. Richtig: $V = 48 \cdot 9 : 3 = 144$ cm³ \quad b) Drei gleiche Pyramiden mit derselben Grundfläche und Höhe füllen den Quader genau; eine Pyramide fasst also ein Drittel. \quad c) Quadratische Pyramide mit Grundkante $4$ cm; die Seitenhöhe $6$ cm läuft auf einer Seitenfläche von der Spitze zur Mitte einer Grundkante \quad d) $M = 4 \cdot 6 \cdot 11 : 2 = 132$ m²; Kosten $132 \cdot 85 = 11\,220$ € – die $11\,000$ € reichen nicht.}

17c)

\pyramide{4}{4}{5.66}{$4$ cm}{}{}

